\documentclass[11pt]{article}
\usepackage[T1]{fontenc}
\usepackage[utf8]{inputenc}
\usepackage{booktabs}
\usepackage{geometry}
\geometry{margin=1in}
\usepackage{amsmath}
\usepackage{array}

\title{Survey Summary: Twenty-Two Methods for Direction-Aware\\ Variant Effect Prediction in nAChRs}
\author{Hari Krishna Mohan Raj}
\date{}

\begin{document}

\maketitle

\begin{abstract}
This note summarises the twenty-two papers retained in the project survey
sheet. It groups them by the technology each one contributes, describes the
problem each is targeting, and maps each paper onto the specific goal of this
project: predicting gain-of-function versus loss-of-function direction for
missense variants in nicotinic acetylcholine receptors (nAChRs).
\end{abstract}

\section{Overview}

The twenty-two papers fall into seven technological clusters, summarised in
Table~\ref{tab:summary}. Six of them predict direction of effect directly and
are the closest comparators to our model; the rest supply either the feature
representations (protein language models, structure prediction, stability
estimation), the benchmark infrastructure, or the structural-biology data
(cryo-EM nAChR structures) that our pipeline builds on.

\begin{table}[ht]
\footnotesize
\caption{Summary of the twenty-two surveyed papers: technology, target, and relevance.}
\label{tab:summary}
\begin{tabular}{@{}p{3.1cm}p{3.3cm}p{3.4cm}p{4.0cm}@{}}
\toprule
\textbf{Paper (year)} & \textbf{Technology} & \textbf{Target} & \textbf{Relevance to our project} \\
\midrule
Brunklaus 2022 (funNCion) & Sequence + topology classifier & Na/Ca channel GOF/LOF & Direct comparator; proves direction-aware prediction works for ion channels \\
Bo{\ss}elmann 2022 & Multi-task SVM & Potassium-channel GOF/LOF & Direct comparator; taxonomy-based multi-task transfer \\
Bo{\ss}elmann 2023 (SCION) & Multi-kernel SVM + HPO terms & Voltage-gated channel effects & Adds clinical phenotype features (HPO) to effect prediction \\
Stein 2023 (LoGoFunc) & Gradient-boosted ensemble & Genome-wide GOF/LOF & Genome-wide baseline; no receptor-specific context (our gap) \\
Zhong 2024 (PreMode) & SE(3)-equivariant GNN & Mode-of-action GOF/LOF & Structure-GNN transfer with few labels \\
Gies 2025 (MissION) & ESM-2 + multi-channel features & Ion-channel GOF/LOF & Closest SOTA; nAChR is not a focal family (our gap) \\
\midrule
Noviello 2021 ($\alpha$7 nAChR) & Cryo-EM + electrophysiology & nAChR gating mechanism & Basis of our open/closed conformational features \\
Mazzaferro 2024 ($\alpha$4$\beta$2 nAChR) & Cryo-EM + voltage clamp & nAChR allostery & Subunit-interface and ligand-site features \\
\midrule
Rives 2021 (ESM-1b) & Masked transformer PLM (650M) & Unsupervised representation & Precursor to the ESM-2 embeddings we use \\
Meier 2021 (ESM-1v) & Zero-shot masked-marginal scoring & Label-free mutational effects & Our zero-shot ESM score; ranks magnitude, not direction \\
Marquet 2022 & ProtT5/ESM-1b embeddings & Embeddings vs.\ conservation & Confirms embeddings carry conservation but no direction \\
Brandes 2023 & ESM-1b + logistic regression & Whole-proteome pathogenicity & Lightweight PLM-head classifier template \\
Lin 2023 (ESMFold) & ESM-2 + folding head & Fast structure prediction & Residue-level structure for under-studied subunits \\
Sledzieski 2024 (DPLM) & LoRA fine-tuning of ESM-2 & Cheap PLM adaptation & Addresses our small-data fine-tuning constraint \\
Chen 2024 (xTrimoPGLM) & 100B unified transformer & Protein understanding/generation & Suggests scale saturates (we saw 650M $\approx$ 150M) \\
Frazer 2021 (EVE) & Bayesian VAE over MSAs & Pathogenicity from evolution & Conservation signal; needs deep MSAs \\
\midrule
Jumper 2021 (AlphaFold2) & Evoformer + structure module & Atomic structure prediction & Supplies structural features and PDB mapping \\
Abramson 2024 (AlphaFold3) & Diffusion model & Complexes, ligands, ions & Ligand-bound binding-site features \\
\midrule
Gao 2023 (PrimateAI-3D) & 3D-CNN on sequence + structure & Clinical variant classification & Structure-aware deep net, transferable to GOF/LOF \\
Blaabjerg 2023 (RaSP) & GNN + ESM-1v embeddings & Stability $\Delta\Delta G$ & Candidate stability feature (Future Work) \\
\midrule
Notin 2023 (ProteinGym) & Standardised DMS benchmark & Fair effect-predictor comparison & Template for a shared nAChR benchmark \\
Cheng 2023 (AlphaMissense) & AlphaFold2 fine-tuned & Proteome-wide pathogenicity & Key pathogenicity comparator; struggles with GOF \\
\bottomrule
\end{tabular}
\end{table}

\section{Technologies in use}

\paragraph{Protein language models (dominant).}
Seven of the twenty-two papers are built on transformer language models trained
on protein sequence alone. Rives et al.\ established the paradigm with ESM-1b;
Meier et al.\ showed a masked model can score mutations zero-shot as a
log-likelihood ratio; and Lin et al.\ repurposed the same representation for
structure prediction (ESMFold). Marquet et al.\ and Brandes et al.\ both
demonstrate that embeddings are useful, cheap features, the former for
conservation and the latter as the input to a logistic-regression pathogenicity
head. Two papers push scale in opposite directions: Chen et al.\ train a
100-billion-parameter model (xTrimoPGLM), while Sledzieski et al.\ use
parameter-efficient fine-tuning (LoRA), which matters most to us because our
dataset is tiny.

\paragraph{Generative and evolutionary models.}
Frazer et al.\ (EVE) train a Bayesian variational autoencoder across multiple
sequence alignments to score pathogenicity from evolution alone. It is the
strongest pure-conservation signal in the survey but requires deep MSAs and,
like all pathogenicity models, does not distinguish GOF from LOF.

\paragraph{Structure prediction.}
AlphaFold2 (Jumper et al.) and AlphaFold3 (Abramson et al.) supply atomic
coordinates: AlphaFold2 for single-chain structure at experimental accuracy, and
AlphaFold3 for complexes that include ligands and ions. Both are directly
relevant because our structural features are extracted from experimental cryo-EM
maps and AlphaFold models.

\paragraph{Structure-aware effect models.}
Gao et al.\ (PrimateAI-3D) use a 3D convolutional network over sequence and
structure, and Blaabjerg et al.\ (RaSP) use a graph neural network over ESM-1v
embeddings to predict stability change ($\Delta\Delta G$). Both show that
structure-aware deep networks extract signal that sequence-only models miss,
which motivates the structural branch of our feature set.

\paragraph{Direction-of-effect predictors.}
Six papers target the exact problem this project addresses. funNCion
(Brunklaus et al.) classifies GOF versus LOF for sodium and calcium channel
variants; two related works by Bo{\ss}elmann et al.\ apply multi-task and
multi-kernel support-vector machines to voltage-gated potassium, sodium, and
calcium channels, the latter adding clinical phenotype (HPO) features; LoGoFunc
(Stein et al.) predicts GOF/LOF genome-wide with a gradient-boosted ensemble;
PreMode (Zhong et al.) uses an SE(3)-equivariant graph neural network; and
MissION (Gies et al.) combines ESM-2 embeddings with structural, residue, and
ontology features in separate channels to reach ROC-AUC 0.925.

\paragraph{Pathogenicity comparator.}
AlphaMissense (Cheng et al.) fine-tunes AlphaFold2 on population-frequency data
to classify the clinical significance of every human missense variant. It is the
key pathogenicity-only baseline in the survey, and its documented difficulty
with gain-of-function variants reinforces the case for direction-aware models.

\paragraph{nAChR structural biology (experimental data).}
Noviello et al.\ solved the human $\alpha$7 receptor in resting, activated, and
desensitised states, and Mazzaferro et al.\ resolved calcium potentiation of
$\alpha$4$\beta$2 in two stoichiometries. These are not machine-learning methods
but supply the state-specific, interface-level structural knowledge our
conformational and structural features are designed to capture.

\section{How the field is targeting the problem}

The survey reveals a clear split in targets. Most papers target
\emph{pathogenicity}: whether a variant is damaging, measured against ClinVar or
deep-mutational-scan labels. A smaller, more recent wave targets
\emph{direction of effect}, motivated by the clinical fact that GOF and LOF
demand opposite treatments. This second wave is where our project sits, and the
six direction-aware papers are its current state of the art.

Critically, none of the direction-aware tools is built for ligand-gated,
pentameric channels specifically. funNCion and the two Bo{\ss}elmann et al.\
methods are limited to voltage-gated channels. LoGoFunc treats every gene
uniformly and has no structural context. PreMode is binary GOF/LOF and
struggles with disordered regions. MissION aggregates diverse channel families
and does not encode subunit identity or pentameric assembly. This is precisely
the gap our model targets: a nAChR-specific, subunit-aware, direction-aware
predictor.

\section{Relevance to our project}

The papers map onto our pipeline at four levels.

\begin{enumerate}
\item \textbf{Direct comparators.} funNCion, the two Bo{\ss}elmann et al.\
models, LoGoFunc, PreMode, and MissION define the baseline we must beat or
match. MissION in particular is the closest precursor, and its channel-agnostic
design is the limitation we exploit by encoding subunit identity and pentameric
context.
\item \textbf{Feature sources.} ESM-2 embeddings and zero-shot scores (Rives,
Meier, Lin, Sledzieski), conservation signals (Frazer, Marquet), structural
features from cryo-EM and AlphaFold models (Jumper, Abramson, Noviello,
Mazzaferro), clinical phenotype terms (Bo{\ss}elmann et al.\ 2023), and
stability estimates (Blaabjerg) are all candidate inputs to our 71-feature
representation.
\item \textbf{Methodological guidance.} Brandes et al.\ show a logistic head on
PLM embeddings is a cheap, effective classifier; Sledzieski et al.\ show LoRA
fine-tuning works with very few trainable parameters, which is exactly the
setting of a 717-variant dataset; Zhong et al.\ show a structure GNN transfers
with few labels; and Gao et al.\ show structure-aware convolution outperforms
sequence-only models.
\item \textbf{Evaluation.} ProteinGym (Notin et al.) is the template for a
shared, versioned benchmark, which our Future Work proposes for nAChR
direction-of-effect prediction.
\end{enumerate}

One cross-cutting observation from the survey is that scaling protein language
models has begun to saturate: Chen et al.\ report only incremental gains at
100B parameters, and our own experiments found ESM-2 650M performs no better
than 150M on this task. The more promising levers for direction-of-effect
prediction are structure, subunit context, and direction-labelled data, not
larger language models.

The twenty-two papers below correspond one-to-one with the survey sheet entries.

\begin{thebibliography}{99}\setlength{\itemsep}{1pt}\small

\bibitem{rives2021} A. Rives et al., Biological structure and function emerge from scaling unsupervised learning to 250 million protein sequences, \emph{PNAS} 118(15), 2021. \texttt{10.1073/pnas.2016239118}.

\bibitem{meier2021} J. Meier et al., Language models enable zero-shot prediction of the effects of mutations on protein function, \emph{bioRxiv}, 2021. \texttt{10.1101/2021.07.09.450648}.

\bibitem{frazer2021} J. Frazer et al., Disease variant prediction with deep protein language models, \emph{Nature Genetics} 53:1512, 2021. \texttt{10.1038/s41588-021-00952-2}.

\bibitem{jumper2021} J. Jumper et al., Highly accurate protein structure prediction with AlphaFold, \emph{Nature} 596:583, 2021. \texttt{10.1038/s41586-021-03819-2}.

\bibitem{noviello2021} C. M. Noviello et al., Structure and gating mechanism of the $\alpha$7 nicotinic acetylcholine receptor, \emph{Cell} 184:2121, 2021. \texttt{10.1016/j.cell.2021.02.049}.

\bibitem{brunklaus2022} A. Brunklaus et al., Gene variant effects across sodium channelopathies predict function and guide precision therapy, \emph{Brain} 145:4275, 2022. \texttt{10.1093/brain/awac006}.

\bibitem{bosselmann2022} C. M. Bo{\ss}elmann et al., Predicting the functional effects of voltage-gated potassium channel missense variants with multi-task learning, \emph{eBioMedicine}, 2022. \texttt{10.1016/j.ebiom.2022.104115}.

\bibitem{marquet2022} C. Marquet et al., Embeddings from protein language models predict conservation and variant effects, \emph{Human Genetics} 141:1629, 2022. \texttt{10.1007/s00439-021-02411-y}.

\bibitem{brandes2023} N. Brandes et al., Genome-wide prediction of disease variant effects with a deep protein language model, \emph{Nature Genetics} 55:1512, 2023. \texttt{10.1038/s41588-023-01465-w}.

\bibitem{lin2023} Z. Lin et al., Evolutionary-scale prediction of atomic-level protein structure with a language model, \emph{Science} 379:1123, 2023. \texttt{10.1126/science.ade2574}.

\bibitem{gao2023} H. Gao et al., The landscape of tolerated genetic variation in humans and primates, \emph{Science} 380:eabn8153, 2023. \texttt{10.1126/science.abn8153}.

\bibitem{stein2023} D. Stein et al., Genome-wide prediction of pathogenic gain- and loss-of-function variants from ensemble learning of a diverse feature set, \emph{Genome Medicine} 15:103, 2023. \texttt{10.1186/s13073-023-01261-9}.

\bibitem{blaabjerg2023} L. M. Blaabjerg et al., Rapid protein stability prediction using deep learning representations, \emph{eLife} 12:e82593, 2023. \texttt{10.7554/eLife.82593}.

\bibitem{notin2023} P. Notin et al., ProteinGym: large-scale benchmarks for protein design and fitness prediction, \emph{bioRxiv}, 2023. \texttt{10.1101/2023.12.07.570727}.

\bibitem{bosselmann2023} C. M. Bo{\ss}elmann et al., Predicting functional effects of ion channel variants using new phenotypic machine learning methods, \emph{PLOS Computational Biology}, 2023. \texttt{10.1371/journal.pcbi.1010959}.

\bibitem{cheng2023} J. Cheng et al., Accurate proteome-wide missense variant effect prediction with AlphaMissense, \emph{Science} 381:eadg7492, 2023. \texttt{10.1126/science.adg7492}.

\bibitem{abramson2024} J. Abramson et al., Accurate structure prediction of biomolecular interactions with AlphaFold3, \emph{Nature} 630:493, 2024. \texttt{10.1038/s41586-024-07487-w}.

\bibitem{zhong2024} G. Zhong et al., PreMode predicts mode-of-action of missense variants by deep graph representation learning of protein sequence and structural context, \emph{bioRxiv}, 2024. \texttt{10.1101/2024.02.20.581321}.

\bibitem{mazzaferro2024} S. Mazzaferro et al., Structural bases for stoichiometry-selective calcium potentiation of a neuronal nicotinic receptor, \emph{British Journal of Pharmacology} 181:1584, 2024. \texttt{10.1111/bph.16321}.

\bibitem{sledzieski2024} S. Sledzieski et al., Democratizing protein language models with parameter-efficient fine-tuning, \emph{PNAS} 121(39), 2024. \texttt{10.1073/pnas.2405840121}.

\bibitem{chen2024} B. Chen et al., xTrimoPGLM: unified 100B-scale pretrained transformer for deciphering the language of protein, \emph{Nature Methods} 21:1506, 2024. \texttt{10.1038/s41592-024-02305-z}.

\bibitem{gies2025} S. Gies et al., Functional effect predictions for ion channel missense variants using a protein language model, \emph{medRxiv}, 2025. \texttt{10.1101/2025.10.16.25337735}.

\end{thebibliography}

\end{document}
